\section{Síntesis y ampliación}

\subsection{Repaso de los puntos clave de esta clase}

Esta clase, como cierre de la serie de clases fundamentales de MIT 6.S191, cubrió tres temas centrales:

\begin{enumerate}
    \item \textbf{La naturaleza y las limitaciones del aprendizaje profundo}
    \begin{itemize}
        \item Las redes neuronales son aproximadores de funciones; el Universal Approximation Theorem garantiza la capacidad teórica, pero no la generalización
        \item El experimento con etiquetas aleatorias reveló la sorprendente capacidad de sobreajuste (memorización) de las redes profundas
        \item Tres modos de falla principales: sesgo de los datos, incertidumbre en escenarios críticos para la seguridad, ataques adversarios
        \item Nueve limitaciones principales: hambre de datos, uso intensivo de cómputo, fragilidad ante ataques adversarios, sesgo algorítmico, mala cuantificación de la incertidumbre, caja negra, necesidad de conocimiento experto, dificultad para codificar estructura, dificultad para extrapolar
    \end{itemize}

    \item \textbf{Nueva frontera I -- Diffusion Models}
    \begin{itemize}
        \item Superan el mode collapse y la inestabilidad de entrenamiento de las VAE/GAN
        \item Construyen pares de entrenamiento agregando ruido hacia adelante, y generan mediante eliminación de ruido hacia atrás
        \item La estrategia de generación iterativa es el núcleo de su ventaja en calidad
        \item Aplicaciones amplias: texto a imagen, generación de video, diseño de moléculas y proteínas
    \end{itemize}

    \item \textbf{Nueva frontera II -- Large Language Models}
    \begin{itemize}
        \item Su naturaleza: redes neuronales a escala masiva más datos de texto a escala masiva
        \item Objetivo de entrenamiento central: Next Token Prediction
        \item Las capacidades emergentes surgen con un crecimiento en forma de «cambio de fase» a medida que aumenta el tamaño del modelo
        \item Retos clave: robustez, alucinaciones, protección de seguridad, razonamiento lógico
    \end{itemize}
\end{enumerate}

\subsection{Lecturas adicionales}

\begin{itemize}
    \item Hornik, K., Stinchcombe, M., \& White, H. (1989). Multilayer feedforward networks are universal approximators. \textit{Neural Networks}.
    \item Zhang, C., et al. (2017). Understanding deep learning requires rethinking generalization. \textit{ICLR 2017}.
    \item Goodfellow, I., Shlens, J., \& Szegedy, C. (2015). Explaining and harnessing adversarial examples. \textit{ICLR 2015}.
    \item Athalye, A., et al. (2018). Synthesizing robust adversarial examples. \textit{ICML 2018}.
    \item Ho, J., Jain, A., \& Abbeel, P. (2020). Denoising diffusion probabilistic models. \textit{NeurIPS 2020}.
    \item Sohl-Dickstein, J., et al. (2015). Deep unsupervised learning using nonequilibrium thermodynamics. \textit{ICML 2015}.
    \item Ramesh, A., et al. (2022). Hierarchical text-conditional image generation with CLIP latents. \textit{arXiv 2022}.
    \item Wei, J., et al. (2022). Emergent abilities of large language models. \textit{TMLR 2022}.
    \item Sitio oficial del curso: \url{http://introtodeeplearning.com}
    \item GitHub: \url{https://github.com/MITDeepLearning/introtodeeplearning/}
\end{itemize}

%% --- Fin del contenido --- %%

\end{document}
